\section{Назначение и границы проекта}
\label{sec:overview}

PharmDelivery моделирует работу службы доставки лекарств в течение заданного
числа дней. Программа генерирует каталог и покупателей, формирует заказы,
выделяет товар со склада, распределяет доставки между курьерами, заказывает
пополнение и рассчитывает финансовый результат.

Модель предназначена для учебных экспериментов: изменения наценки, количества
курьеров, ассортимента, скидок и начального запаса. Данные покупателей и каталога
генерируются внутри программы; входной CSV позволяет задать начальные партии
склада. Итоги доступны в консоли, через структуры C++ и в JSON-файле.

Время дискретно: день эксперимента имеет номер от 1 до $N$, а день 0 обозначает
инициализацию. Денежные значения измеряются в рублях, количества --- в упаковках,
периоды и сроки годности --- в днях. Все идентификаторы лекарств и индексы
курьеров в программных структурах начинаются с нуля.

\subsection{Основные обозначения}
\begin{tabularx}{\textwidth}{@{}L{0.18\textwidth}Y@{}}
\toprule
Обозначение & Значение \\
\midrule
$N$ & Продолжительность эксперимента: от 10 до 25 дней. \\
$M$ & Количество курьеров: от 3 до 9. \\
$K$ & Количество лекарств в каталоге: от 15 до 35. \\
$m$ & Наценка в процентах; значение 25 означает 25\%. \\
$s_i$ & Минимальный остаток лекарства $i$: поле \code{min_stock}. \\
$e$ & Последний день годности партии: поле \code{expiration_day}. \\
\bottomrule
\end{tabularx}

\subsection{Реализованные возможности}
\begin{itemize}
  \item Воспроизводимая генерация каталога, заказов и поставок по \code{seed}.
  \item Постоянные покупатели с периодическими покупками и поток разовых заказов.
  \item Полная и частичная выдача из нескольких партий; контроль срока годности
        на предполагаемый день доставки.
  \item Однократная уценка близких к истечению партий и списание просроченных.
  \item Скидки по карте, за крупную покупку и для постоянного клиента.
  \item Пополнение дефицитных лекарств с задержкой от одного до трёх дней.
  \item Очередь при нехватке вместимости доставки, учёт недогрузки и опозданий.
  \item Контроль финансового баланса и подробный JSON-отчёт.
\end{itemize}

\subsection{Границы реализации}
В текущем проекте отсутствуют HTTP-сервер, REST-маршруты, база данных,
авторизация и сетевые интеграции. Термин backend здесь обозначает вычислительную
модель и её программный интерфейс. Qt нужен для консольной точки входа
с \code{QApplication}; окно интерфейса и цикл \code{exec()} не создаются.

Курьер представлен индексом в массиве нагрузки; отдельного класса курьера,
координат, маршрутов и времени поездки нет. Доставка занимает один шаг после
подготовки заказа. Расходы на транспорт, оплату труда и хранение в финансовую
модель не включены.

Глобальные данные и единственный экземпляр движка рассчитаны на один эксперимент
в процессе. Параллельный запуск нескольких независимых моделей через этот
контракт не реализован. Повторный последовательный эксперимент выполняется
через новую настройку и инициализацию.
